% =====================================================================
% §6 NSGA-II — ANÁLISE ISOLADA
% Lida separadamente por causa do caveat de factibilidade (feedback Maristela)
% =====================================================================
\section{NSGA-II: Análise Isolada}

%%% Frame 6.1 — Factibilidade (o caveat, dito com clareza)
\begin{frame}{Factibilidade das soluções NSGA-II}

    \begingroup
    \setbeamercolor{block title}{fg=white, bg=red!40}
    \setbeamercolor{block body}{bg=red!10}
    \begin{block}{Por que analisar separadamente}
        Nenhuma solução NSGA-II satisfaz a cobertura mínima (R\ref{eq:r6}). Sem factibilidade garantida, \textbf{não há comparação direta} com métodos que a garantem (Exato, MIP-heurísticas).
    \end{block}
    \endgroup

    \begin{block}{Distinção importante}
        \begin{itemize}
            \item Restrições estruturais R\ref{eq:r1}--R\ref{eq:r4}: satisfeitas (Estendido 100\%; Compacto 89,8\%).
            \item Violação concentrada: cobertura mínima (R\ref{eq:r6}), 4--6\% dos períodos, turno vespertino.
            \item Compacto: 10,2\% também violam capacidade de sala (R\ref{eq:r5}), clínica U06.
        \end{itemize}
    \end{block}
    % TODO: valores de Z1/Z2 verificados via Gurobi com variáveis fixadas.
\end{frame}

%%% Frame 6.2 — Discussão isolada (Compacto >> Estendido)
\begin{frame}{Discussão isolada: Compacto \emph{vs.} Estendido}

    \begin{block}{Dentro da família meta-heurística}
        \begin{itemize}
            \item Compacto executa muito mais gerações (cromossomo $\sim$70$\times$ menor).
            \item Compacto \textbf{domina} o Estendido em todos os horizontes.
            \item Escala para 60--90 dias produzindo fronteiras densas ($\sim$100 soluções).
        \end{itemize}
    \end{block}

    \begingroup
    \setbeamercolor{block title}{fg=white, bg=UFTblue}
    \setbeamercolor{block body}{bg=UFTblue!10}
    \begin{block}{Encaminhamento}
        Limitação de factibilidade $\rightarrow$ \textbf{operador de reparo} (trabalho futuro), não comparação direta com o exato.
    \end{block}
    \endgroup
    % TODO: leitura sempre isolada, dentro da família NSGA-II.
\end{frame}
